\documentclass[12pt,a4paper]{article}

\usepackage[margin=1.4cm]{geometry}
\usepackage{array}

\usepackage{enumitem}

% Global compact spacing for lists
\setlist[itemize]{
	topsep=0pt,
	partopsep=0pt,
	parsep=0pt,
	itemsep=0pt
}

\setlist[enumerate]{
	topsep=0pt,
	partopsep=0pt,
	parsep=0pt,
	itemsep=0pt
}

\usepackage{multicol}

\newcolumntype{L}[1]{>{\raggedright\arraybackslash}p{#1}}

\newcommand{\checkbox}{\fbox{\rule{0pt}{0.8ex}\rule{0.8ex}{0pt}}}

\begin{document}
	

	\begin{center}
		{\Large\bfseries C1: OPORTUNIDADES A NUESTRO ALREDEDOR}
		
		\vspace{0.1cm}
		
		Emprendimiento e Innovación Grado 10
		
		\vspace{0.3cm}
		
		\begin{tabular}{L{0.13\textwidth}L{0.65\textwidth}}
			Grado: & \rule{2cm}{0.4pt} \\
			Fecha: & \rule{2cm}{0.4pt} \\
			Nombre 1: & \rule{\linewidth}{0.4pt} \\
			Nombre 2: & \rule{\linewidth}{0.4pt} \\
			Nombre 3: & \rule{\linewidth}{0.4pt} \\
		\end{tabular}
	\end{center}
	
	\section*{Nuestro objetivo}
	
	La innovación comienza cuando identificamos algo que podría mejorarse. Primero, generarán ideas a partir de su experiencia en el aula. Luego, como tarea, observarán su entorno e identificarán situaciones que puedan representar una necesidad, un problema, una incomodidad o una oportunidad de mejora.
	
	En esta etapa, no diseñen una solución. El objetivo es identificar situaciones que valga la pena investigar.
	
	\section*{Actividad en clase: Generar ideas}
	
	Piensen en su experiencia en el colegio y en el aula. Consideren situaciones que:
	
	\begin{tabular}{@{}p{0.43\linewidth}@{\hspace{0.04\linewidth}}p{0.53\linewidth}@{}}
		\begin{itemize}[leftmargin=0.5cm,itemsep=0pt]
			\item Hacen que una actividad sea difícil o poco eficiente.
			\item Generan incomodidad o frustración.
			\item Desperdician tiempo, materiales o recursos.
		\end{itemize}
		&
		\begin{itemize}[leftmargin=0.5cm,itemsep=0pt]
			\item Generan dificultades para estudiantes, profesores u otras personas.
			\item Podrían organizarse o mejorarse.
			\item Podrían realizarse de una manera diferente.
		\end{itemize}
	\end{tabular}
	
	\vspace{0.3cm}
	
	\textbf{Ideas de nuestro salón de clase:}
	
	\vspace{0.1cm}
	
	1.\ \rule{0.85\linewidth}{0.4pt}
	
	\vspace{0.25cm}
	
	2.\ \rule{0.85\linewidth}{0.4pt}
	
	\vspace{0.25cm}
	
	3.\ \rule{0.85\linewidth}{0.4pt}
	
	\vspace{0.25cm}
	
	4.\ \rule{0.85\linewidth}{0.4pt}
	
	\vspace{-5mm}
	\section*{Observen su entorno}
	
	Durante su día normal, presten atención a lo que sucede a su alrededor. Pueden observar:
	
	\begin{multicols}{2}
		\begin{itemize}[leftmargin=0.6cm,itemsep=0pt]
			\item Su colegio
			\item Su hogar
			\item Su barrio
			\item Espacios públicos
			\item Transporte
			\item Tiendas o negocios
			\item Espacios recreativos
			\item Situaciones relacionadas con la tecnología
			\item Otros entornos de la comunidad.
		\end{itemize}
	\end{multicols}
	
	Busquen situaciones en las que las personas experimenten una dificultad, necesidad, incomodidad, ineficiencia o una oportunidad que no esté siendo atendida. \vspace{6mm}
	
	\begin{tabular}{@{}p{0.32\linewidth}p{0.63\linewidth}@{}}
		No comiencen preguntándose:
		&
		\textit{``¿Qué debería construir?''}
		\\[0.15cm]
		&\\
		En cambio, pregúntense:
		&
		\textit{``¿Qué está sucediendo y qué podría mejorarse?''}
	\end{tabular}
	
	\newpage
	
	\section*{Qué deben registrar}
	
	Para cada oportunidad, registren:
	
	\begin{center}
		\small
		\begin{tabular}{|L{0.20\textwidth}|L{0.26\textwidth}|L{0.38\textwidth}|}
			\hline
			\textbf{Elemento} & \textbf{Qué escribir} & \textbf{Ejemplo} \\
			\hline
			\textbf{Lugar} & ¿Dónde? & Cafetería del colegio \\
			\hline
			\textbf{Situación} & ¿Qué está sucediendo? & Los estudiantes esperan en una fila larga. \\
			\hline
			\textbf{Necesidad / Problema / Incomodidad} & ¿Qué necesita mejorar? & Los estudiantes pasan tiempo esperando. \\
			\hline
			\textbf{¿A quién afecta?} & ¿Quién lo experimenta? & Estudiantes y personal de la cafetería. \\
			\hline
			\textbf{Evidencia} & ¿Qué observaron? & La fila se hizo muy larga. \\
			\hline
			\textbf{¿Por qué importa?} & ¿Por qué vale la pena investigarlo? & Los estudiantes tienen menos tiempo para comer. \\
			\hline
			\textbf{Posible solución} & ¿Ya tienen una? & Dejar abierta por ahora. \\
			\hline
		\end{tabular}
	\end{center}
	
	\textbf{Importante:} Primero comprendan la situación. No decidan la solución demasiado pronto.
	
	\section*{Oportunidad 1}
	
	\begin{center}
		\small
		\begin{tabular}{|L{0.22\textwidth}|L{0.61\textwidth}|}
			\hline
			\textbf{Lugar} & \rule{0pt}{0.5cm} \\
			\hline
			\textbf{Situación} & \rule{0pt}{0.5cm} \\
			\hline
			\textbf{Necesidad / Problema / Incomodidad} & \rule{0pt}{0.5cm} \\
			\hline
			\textbf{¿A quién afecta?} & \rule{0pt}{0.45cm} \\
			\hline
			\textbf{Evidencia} & \rule{0pt}{0.6cm} \\
			\hline
			\textbf{¿Por qué importa?} & \rule{0pt}{0.5cm} \\
			\hline
			\textbf{Posible solución} & \rule{0pt}{0.45cm} \\
			\hline
		\end{tabular}
	\end{center}
	
	\vspace{0.15cm}
	
	\section*{Oportunidad 2}
	
	\begin{center}
		\small
		\begin{tabular}{|L{0.22\textwidth}|L{0.61\textwidth}|}
			\hline
			\textbf{Lugar} & \rule{0pt}{0.5cm} \\
			\hline
			\textbf{Situación} & \rule{0pt}{0.5cm} \\
			\hline
			\textbf{Necesidad / Problema / Incomodidad} & \rule{0pt}{0.5cm} \\
			\hline
			\textbf{¿A quién afecta?} & \rule{0pt}{0.45cm} \\
			\hline
			\textbf{Evidencia} & \rule{0pt}{0.6cm} \\
			\hline
			\textbf{¿Por qué importa?} & \rule{0pt}{0.5cm} \\
			\hline
			\textbf{Posible solución} & \rule{0pt}{0.45cm} \\
			\hline
		\end{tabular}
	\end{center}
	
	\vspace{0.15cm}
	
	\section*{Oportunidad 3}
	
	\begin{center}
		\small
		\begin{tabular}{|L{0.22\textwidth}|L{0.61\textwidth}|}
			\hline
			\textbf{Lugar} & \rule{0pt}{0.5cm} \\
			\hline
			\textbf{Situación} & \rule{0pt}{0.5cm} \\
			\hline
			\textbf{Necesidad / Problema / Incomodidad} & \rule{0pt}{0.5cm} \\
			\hline
			\textbf{¿A quién afecta?} & \rule{0pt}{0.45cm} \\
			\hline
			\textbf{Evidencia} & \rule{0pt}{0.6cm} \\
			\hline
			\textbf{¿Por qué importa?} & \rule{0pt}{0.5cm} \\
			\hline
			\textbf{Posible solución} & \rule{0pt}{0.45cm} \\
			\hline
		\end{tabular}
	\end{center}
	
	\section*{Oportunidad 4}
	
	\begin{center}
		\small
		\begin{tabular}{|L{0.22\textwidth}|L{0.61\textwidth}|}
			\hline
			\textbf{Lugar} & \rule{0pt}{0.5cm} \\
			\hline
			\textbf{Situación} & \rule{0pt}{0.5cm} \\
			\hline
			\textbf{Necesidad / Problema / Incomodidad} & \rule{0pt}{0.5cm} \\
			\hline
			\textbf{¿A quién afecta?} & \rule{0pt}{0.45cm} \\
			\hline
			\textbf{Evidencia} & \rule{0pt}{0.6cm} \\
			\hline
			\textbf{¿Por qué importa?} & \rule{0pt}{0.5cm} \\
			\hline
			\textbf{Posible solución} & \rule{0pt}{0.45cm} \\
			\hline
		\end{tabular}
	\end{center}
	
	\vspace{0.15cm}
	
	\section*{Oportunidad 5}
	
	\begin{center}
		\small
		\begin{tabular}{|L{0.22\textwidth}|L{0.61\textwidth}|}
			\hline
			\textbf{Lugar} & \rule{0pt}{0.5cm} \\
			\hline
			\textbf{Situación} & \rule{0pt}{0.5cm} \\
			\hline
			\textbf{Necesidad / Problema / Incomodidad} & \rule{0pt}{0.5cm} \\
			\hline
			\textbf{¿A quién afecta?} & \rule{0pt}{0.45cm} \\
			\hline
			\textbf{Evidencia} & \rule{0pt}{0.6cm} \\
			\hline
			\textbf{¿Por qué importa?} & \rule{0pt}{0.5cm} \\
			\hline
			\textbf{Posible solución} & \rule{0pt}{0.45cm} \\
			\hline
		\end{tabular}
	\end{center}
	
	\vspace{0.15cm}
	
	\section*{Oportunidad 6}
	
	\begin{center}
		\small
		\begin{tabular}{|L{0.22\textwidth}|L{0.61\textwidth}|}
			\hline
			\textbf{Lugar} & \rule{0pt}{0.5cm} \\
			\hline
			\textbf{Situación} & \rule{0pt}{0.5cm} \\
			\hline
			\textbf{Necesidad / Problema / Incomodidad} & \rule{0pt}{0.5cm} \\
			\hline
			\textbf{¿A quién afecta?} & \rule{0pt}{0.45cm} \\
			\hline
			\textbf{Evidencia} & \rule{0pt}{0.6cm} \\
			\hline
			\textbf{¿Por qué importa?} & \rule{0pt}{0.5cm} \\
			\hline
			\textbf{Posible solución} & \rule{0pt}{0.45cm} \\
			\hline
		\end{tabular}
	\end{center}
	
	\section*{Reflexión inicial}
	
	\textbf{1. ¿Qué oportunidad parece más interesante para investigar? ¿Por qué?}
	
	\vspace{0.1cm}
	
	\rule{\linewidth}{0.4pt}
	
	\vspace{0.35cm}
	
	\rule{\linewidth}{0.4pt}
	
	\vspace{0.4cm}
	
	\textbf{2. ¿Qué evidencia les hace pensar que esta es una situación real y no una suposición?}
	
	\vspace{0.1cm}
	
	\rule{\linewidth}{0.4pt}
	
	\vspace{0.35cm}
	
	\rule{\linewidth}{0.4pt}
	
	\vspace{0.4cm}
	
	\textbf{3. ¿Quién podría ayudarles a comprender mejor esta situación?}
	
	\vspace{0.1cm}
	
	\rule{\linewidth}{0.4pt}
	
	\vspace{0.35cm}
	
	\rule{\linewidth}{0.4pt}
	
	\vspace{0.4cm}
	
	\textbf{4. ¿Qué necesitarían averiguar todavía?}
	
	\vspace{0.1cm}
	
	\rule{\linewidth}{0.4pt}
	
	\vspace{0.35cm}
	
	\rule{\linewidth}{0.4pt}
	
	\section*{Verificación final}
	
	\noindent\checkbox\quad Identifiqué situaciones de mi entorno real.
	
	\vspace{0.15cm}
	
	\noindent\checkbox\quad Describí lo que realmente observé.
	
	\vspace{0.15cm}
	
	\noindent\checkbox\quad Identifiqué quiénes están afectados.
	
	\vspace{0.15cm}
	
	\noindent\checkbox\quad Incluí evidencia en lugar de basarme únicamente en opiniones.
	
	\vspace{0.15cm}
	
	\noindent\checkbox\quad Expliqué por qué la situación podría ser importante.
	
	\vspace{0.15cm}
	
	\noindent\checkbox\quad No asumí una solución final demasiado pronto.
	
	\vspace{0.3cm}
	
	\begin{center}
		\textbf{Primero comprendan la situación. Luego comprendan a las personas.}
	\end{center}

	
\end{document}
